% ============================================================
% 中文报告前导（配合 article / extarticle + xeCJK）
%
% 用法：
%   pandoc src/report.qmd -o out.pdf --pdf-engine=xelatex \
%     --include-in-header=templates/header-cn.tex ...
%   -V documentclass=article -V fontsize=11pt -V papersize=a4 \
%   -V "geometry:top=2.2cm,bottom=2.2cm,left=2.5cm,right=2.5cm"
%
% 要改页眉文字，在项目里另建一个 wrapper：
%   \newcommand{\doctitle}{……} \newcommand{\docauthor}{……}
%   \input{<本文件>}
% 或者直接在下面改默认值。
%
% 纯中文文档若有 ctex 可用，优先用 \documentclass{ctexart}——它自带 \zihao、
% 断行与中文图表名，不需要这一份手工前导。
% ============================================================

% ---- 字体：正文 Times + 宋体，标题黑体，等宽 Consolas ----
\usepackage{xeCJK}
\setCJKmainfont[AutoFakeBold=2.2]{SimSun}
\setCJKsansfont[AutoFakeBold=2.2]{SimHei}
\setCJKmonofont{SimSun}
\setmainfont{Times New Roman}
\setsansfont{Arial}
\setmonofont{Consolas}

% ---- 段落：中文习惯——首行缩进 2 字符、段间距 0 ----
\linespread{1.17}              % ≈1.4× 字高；正文 11pt → 基线距 15.7pt
\setlength{\parindent}{2em}
\setlength{\parskip}{0pt}
\setlength{\emergencystretch}{2em}
\widowpenalty=10000
\clubpenalty=10000
\raggedbottom

% ---- 标题：黑体，逐级递减 ----
\usepackage{titlesec}
\titleformat{\section}{\sffamily\fontsize{14pt}{18pt}\selectfont\bfseries}{\thesection}{0.7em}{}
\titlespacing*{\section}{0pt}{15pt plus 3pt minus 2pt}{7pt plus 1pt}
\titleformat{\subsection}{\sffamily\fontsize{12pt}{15pt}\selectfont\bfseries}{\thesubsection}{0.6em}{}
\titlespacing*{\subsection}{0pt}{12pt plus 2pt minus 2pt}{5pt plus 1pt}
\titleformat{\subsubsection}{\sffamily\fontsize{11pt}{14pt}\selectfont\bfseries}{\thesubsubsection}{0.6em}{}
\titlespacing*{\subsubsection}{0pt}{10pt plus 2pt minus 2pt}{4pt}

% ---- 交叉引用宏 ----
% pandoc 的 markdown 里直接写 "表~\ref{...}"，那个 ~ 会被转成字面的波浪号；
% 用下面两个宏把中文前缀和不可断空格一起封进 raw LaTeX。
\newcommand{\tabref}[1]{表~\ref{#1}}
\newcommand{\figref}[1]{图~\ref{#1}}

% ---- 图表题注：表题在上、图题在下，编号后留 1em ----
\usepackage[font=small,labelsep=quad,skip=6pt,justification=centering]{caption}
\captionsetup[table]{position=top}
\renewcommand{\tablename}{表}
\renewcommand{\figurename}{图}
\renewcommand{\thetable}{\arabic{table}}
\renewcommand{\thefigure}{\arabic{figure}}

% ---- 三线表：中文规范的 1.5 / 0.75 / 1.5pt（booktabs 默认只有 0.8/0.5pt）----
\usepackage{booktabs}
\usepackage{tabularx}
\usepackage{array}
\setlength{\heavyrulewidth}{1.5pt}
\setlength{\lightrulewidth}{0.75pt}
\setlength{\cmidrulewidth}{0.5pt}
\renewcommand{\arraystretch}{1.10}
\newcolumntype{L}[1]{>{\raggedright\arraybackslash}p{#1}}

% ---- 图与浮动体 ----
\usepackage{graphicx}
\usepackage{float}
\setlength{\textfloatsep}{14pt plus 3pt minus 3pt}
\setlength{\intextsep}{8pt plus 2pt minus 2pt}
\renewcommand{\topfraction}{0.92}
\renewcommand{\bottomfraction}{0.80}
\renewcommand{\textfraction}{0.08}
\renewcommand{\floatpagefraction}{0.75}

% ---- 页眉页脚 ----
\usepackage{fancyhdr}
\usepackage{etoolbox}
\providecommand{\doctitle}{}
\providecommand{\docauthor}{}
% 可选：页眉右侧在作者后面再接一段（课程名、指导教师之类），默认空
\providecommand{\docrightextra}{}
\pagestyle{fancy}
\fancyhf{}
\fancyhead[L]{\footnotesize\sffamily\doctitle}
\fancyhead[R]{\footnotesize\sffamily\docauthor%
  \ifstrempty{\docrightextra}{}{\quad\docrightextra}}
\fancyfoot[C]{\footnotesize\thepage}
\renewcommand{\headrulewidth}{0.4pt}
\fancypagestyle{plain}{\fancyhf{}\fancyfoot[C]{\footnotesize\thepage}%
  \renewcommand{\headrulewidth}{0pt}}
\setlength{\headheight}{15pt}

% ---- 目录 ----
% pandoc --toc 会插 \tableofcontents；标签默认是英文，这里改成中文。
\renewcommand{\contentsname}{目录}

% ---- 参考文献：默认 citeproc + nature.csl ----
% 要走 GB/T 7714（中文课程报告/学位论文），build.py 的 citations: gbt 会先 include
% templates/gbt-on.tex 定义 \usegbt，这里再加载 gbt7714；样式由 -V biblio-style 给。
% 注意：不要在头文件里写 \bibliographystyle——pandoc 模板已按 biblio-style 输出一条，
% 两条会让 bibtex 报 "Illegal, another \bibstyle command" 并放弃整个文献表。
\ifdefined\usegbt
  \usepackage{gbt7714}
\fi

% ---- 链接与 PDF 元数据 ----
\usepackage[unicode,hidelinks]{hyperref}
\AtBeginDocument{\hypersetup{pdftitle={\doctitle},pdfauthor={\docauthor}}}
